\section{Luồng và lát cắt}
\begingroup\small
\begin{multicols}{2}
\subsection{Mô hình mạng luồng}
Mạng có nguồn $s$, đích $t$, mỗi cạnh $e=(u,v)$ có sức chứa
$c_e\ge0$. Luồng $f_e$ thỏa $0\le f_e\le c_e$ và tại mọi đỉnh
khác $s,t$, tổng luồng vào bằng tổng luồng ra.
Giá trị luồng là luồng ròng ra khỏi $s$.
Đồ thị thặng dư chứa cạnh thuận sức chứa $c_e-f_e$ và cạnh
ngược sức chứa $f_e$; đi cạnh ngược nghĩa là hủy bớt luồng
đã gửi. Mỗi cạnh được thêm vào code tạo một cặp cạnh thặng dư.

\subsection{Luồng cực đại và lát cắt cực tiểu}
Lát cắt $(S,V\setminus S)$ thỏa $s\in S,t\notin S$ có sức chứa
$\sum_{u\in S,v\notin S}c(u,v)$. Mọi luồng có giá trị không
vượt sức chứa mọi lát cắt. Khi không còn đường tăng từ $s$
tới $t$ trong đồ thị thặng dư, tập đỉnh còn tới được từ $s$
cho một lát cắt có sức chứa đúng bằng giá trị luồng;
đó là định lý max-flow min-cut.
Nếu mọi sức chứa nguyên, tồn tại luồng cực đại nguyên.

\subsection{Thuật toán Dinic}
Lặp: BFS trên cạnh thặng dư dương để tạo đồ thị mức;
DFS chỉ đi cạnh làm mức tăng $1$ để đẩy một luồng chặn.
Khi không tới được $t$, dừng. Mỗi pha xóa các đường ngắn nhất
trong đồ thị thặng dư; số pha không quá $V-1$, độ phức tạp tổng
quát $O(V^2E)$. Mẫu code nhận đỉnh $[0,n)$, yêu cầu $s\ne t$;
\texttt{maxflow} trả lượng luồng đẩy thêm trong lần gọi hiện tại,
\texttt{minCutSide} trả các đỉnh ở phía nguồn của lát cắt tối tiểu
sau khi chạy luồng cực đại.
Tổng luồng phải vừa \texttt{long long}; DFS đệ quy có thể sâu $O(V)$.

\paragraph{Mẫu Dinic (C++17).}
\cppsnippet{code/flow/dinic.hpp}

\subsection{Các mẫu dựng mạng từ đề bài}
Đường đi không chung cạnh: đặt sức chứa $1$ trên mỗi cạnh.
Không chung đỉnh: tách $v$ thành $v_{\rm in}\to v_{\rm out}$
sức chứa $1$, trừ nguồn/đích nếu được dùng nhiều lần.
Ghép cặp hai phía: $s\to L$, $L\to R$, $R\to t$, mọi cạnh
sức chứa $1$. Chọn dự án có lãi $w_i$: nếu $w_i>0$, nối
$s\to i$ sức chứa $w_i$; nếu $w_i<0$, nối $i\to t$
sức chứa $-w_i$; điều kiện chọn $i$ buộc chọn $j$ tạo cạnh
$i\to j$ sức chứa \emph{lớn hơn tổng mọi lãi dương}.
Đáp án là tổng lãi dương trừ sức chứa lát cắt tối thiểu.
Tính ngưỡng vô cực theo dữ liệu, tránh tràn số.

\subsection{Luồng có cận dưới}
Với $l_e\le f_e\le u_e$, đặt $f_e=l_e+f'_e$,
$0\le f'_e\le u_e-l_e$. Với mỗi cạnh $u\to v$, cộng
$d_v\mathrel{+}=l_e$, $d_u\mathrel{-}=l_e$.
Thêm siêu nguồn $S$: nếu $d_v>0$, thêm $S\to v$ sức chứa
$d_v$; nếu $d_v<0$, thêm $v\to T$ sức chứa $-d_v$.
Tồn tại lưu thông hợp lệ khi và chỉ khi luồng cực đại từ
$S$ tới $T$ bão hòa toàn bộ cạnh đi ra $S$.
Để tìm luồng $s$--$t$ không âm, thêm cạnh phụ $t\to s$ sức chứa
đủ lớn trước khi kiểm tra. Nếu khả thi, gọi $F_0$ là luồng trên
cạnh phụ. Muốn tối đa hóa: xóa các cạnh nối siêu nguồn/đích
và cạnh phụ $t\to s$, \emph{kể cả các cạnh ngược tương ứng};
giữ các cạnh dư của mạng gốc rồi đẩy thêm luồng từ $s$ tới $t$.
Đáp án là $F_0$ cộng lượng đẩy thêm, và $f_e=l_e+f'_e$.

\subsection{Luồng chi phí cực tiểu}
Mỗi cạnh còn có chi phí $w_e$ cho một đơn vị luồng; cạnh
ngược trong đồ thị thặng dư có chi phí $-w_e$.
Để tìm luồng $F$ đơn vị có chi phí nhỏ nhất, lặp đường ngắn
nhất $s\to t$ trong đồ thị thặng dư và đẩy theo sức chứa nút cổ
chai; dừng khi đủ $F$ hoặc không còn đường. Với thế đỉnh
$\pi$, chi phí giảm của cạnh là
$w'(u,v)=w(u,v)+\pi_u-\pi_v$. Sau khi khởi tạo thế bằng
Bellman--Ford, các chi phí giảm trên phần tới được không âm,
nên có thể dùng Dijkstra mỗi pha.
Mẫu dưới khởi tạo Bellman--Ford từ một siêu nguồn ngầm nối tới
mọi đỉnh, và từ chối mạng dư ban đầu có chu trình chi phí âm,
kể cả ở thành phần rời nguồn. Chu trình âm có sức chứa hữu hạn
không đồng nghĩa đáp án $-\infty$, nhưng cần thuật toán xử lý
lưu thông âm riêng. Mẫu code nhận đỉnh $[0,n)$, yêu cầu
$s\ne t$, lượng cần gửi không âm.
Hàm trả lượng đẩy thêm và chi phí của lượng đó; khi không đủ
đường đi, lượng trả về có thể nhỏ hơn yêu cầu.
Chi phí, khoảng cách và thế đỉnh dùng \texttt{\_\_int128};
tổng chi phí phải vừa kiểu này. Với $A$ lần tăng luồng,
thời gian là $O(VE+A E\log V)$.
Số pha có thể lớn; khi $F$ rất lớn, cần phân tích giới hạn
cụ thể trước khi chọn thuật toán này.

\paragraph{Mẫu min-cost flow (C++17).}
\cppsnippet{code/flow/min_cost_flow.hpp}
\end{multicols}
\endgroup
